\section{Enumeration Protocol}
\label{sec:enumeration-protocol}

This appendix summarizes the finite search spaces used to produce the witness
and certificate artifacts cited in the main text. The guiding rule throughout
the supplementary material is that finite enumeration can support reproducible evidence, but
it does not upgrade a claim to theorem status unless paired with a formal
argument.

\subsection[1-NN witness search]{$1$-NN witness search}

The baseline search enumerates connected simple undirected graphs up to the
configured vertex cutoff, all binary labelings, and all sample orderings under
the deterministic conventions of \Cref{sec:definitions}. For each candidate
sample it computes the fixed-sample maxima of the LOO and replace-one
indicators and records witnesses where the former is $0$ and the latter is $1$.

\subsection{Tie-free filtering}

The tie-free post-processing stage rechecks accepted witness records and keeps
only those whose decisive predictions involve a unique nearest sample occurrence
in the original and perturbed configurations. In the current artifact set, the
minimal two-vertex witness survives this filter.

\subsection[Odd-k gadget candidate search]{Odd-$k$ gadget candidate search}

The odd-$k$ search uses bounded deterministic constructions designed around
repeated-occurrence margin amplification. The output is summarized in
\Cref{tab:k-gadget-candidates} and is explicitly marked as computational
evidence only.

\input{outputs/tables/k_gadget_candidates.tex}

\input{outputs/tables/reproducibility_summary.tex}
